\section{Discussion}
\label{sec:discussion}

\subsection{Jaggedness Tracks Mean Wiggle Differently at Different Pressure Levels}
\label{sec:discussion-jagged}

If we define a judge's \textbf{jaggedness} as the standard deviation of its mean wiggle rates across all datasets, we can plot each judge's mean wiggle against its jaggedness, separately at each pressure level (Figure~\ref{fig:wiggliness-jaggedness}).

\begin{figure}[t]
  \centering
  \includegraphics[width=\linewidth]{per_level_scatter_combined.pdf}
  \caption{Each panel plots all 9 judges at one pressure level: x-axis is the judge's mean wiggle rate at that level (averaged across the 14 judging tasks); y-axis is its cross-dataset jaggedness (std of wiggle rates). Dashed line is the OLS fit, $R^2$ and Pearson $r$ annotated.}
  \label{fig:wiggliness-jaggedness}
\end{figure}

At L1--L3 (mild doubt, counter-argument, expert authority), mean wiggle and jaggedness are strongly positively correlated ($R^2 = 0.68$, $0.94$, $0.87$). When mean wiggle is near zero, std is bounded near zero too, so a low-wiggle judge has nowhere to be jagged. At L4 (consensus pressure) and L5 (strategy cycling), the correlation persists at $R^2 = 0.58$ and $0.75$ but is no longer floor-bounded. Judges that wiggle more on average also have a wider spread of wiggle rates across datasets. At L6 (adaptive persuader), the relationship \emph{inverts} ($R^2 = 0.64$, $r = -0.80$). The judges with the lowest mean wiggle rates (Gemini 3 Flash, Grok-4.1 R, Gemini 3.1 Pro) have \emph{more} cross-dataset spread. Jaggedness (defined as the standard deviation of mean wiggle rates) itself is jagged across different types of pressure.

\subsection{Jury Majority Strength at Baseline Is a Simple Reliability Screen for Golden Sets}
\label{sec:discussion-jury}

If wiggle is jagged, some items are more epistemically uncertain than others, and so we ask: Can a small test predict which? We compare three candidate predictors of per-item wiggle rate: \emph{jury majority strength} (size of the L0 majority across the 9 judges, no pressure applied), \emph{repeat consistency} (temperature-zero per-item agreement), and \emph{position invariance} (verdict survival under argument reordering). Jury majority strength wins at every level (Table~\ref{tab:predictor-comparison}, mean $|\rho| = 0.59$ vs.\ 0.42 for repeat and 0.37 for invariance) and is universal: all 84 (dataset, rubric, scale, level) cells show negative $\rho$, with median $|\rho| = 0.58$ (per-cell table in Appendix~\ref{sec:appendix-jury-heatmap}).

\begin{table}[t]
\centering
\small
\caption{Mean $|\rho|$ between each predictor and per-item wiggle rate, by level. Jury averaged over 84 (dataset, rubric, scale, level) cells; Repeat and Invariance over 72 (not measured on WildGuard). Per-cell table in Appendix~\ref{sec:appendix-jury-heatmap}.}
\label{tab:predictor-comparison}
\begin{tabular}{@{}lrrrrrrr@{}}
\toprule
\textbf{Predictor} & \textbf{L1} & \textbf{L2} & \textbf{L3} & \textbf{L4} & \textbf{L5} & \textbf{L6} & \textbf{Overall} \\
\midrule
Jury Majority Strength & \textbf{0.65} & \textbf{0.57} & \textbf{0.57} & \textbf{0.60} & \textbf{0.59} & \textbf{0.57} & \textbf{0.59} \\
Repeat (temp=0)        & 0.44 & 0.38 & 0.39 & 0.42 & 0.42 & 0.44 & 0.42 \\
Invariance (position)  & 0.35 & 0.36 & 0.38 & 0.35 & 0.37 & 0.38 & 0.37 \\
\bottomrule
\end{tabular}
\end{table}

Beyond predicting wiggle, baseline jury majority strength is \emph{interpretable}: items where 9 frontier models cannot agree at L0 sit in genuinely contested or ambiguous regions of the decision boundary, and more likely regularized over any individual model's idiosyncrasies. A low-majority item is a signal that the underlying question is hard to label --- making it both a wiggle predictor and a flag for content that may not deserve a confident gold label in the first place. Mechanical probes still capture a meaningful fraction of the per-item fragility (mean $|\rho| = 0.42$ for repeat, 0.37 for position invariance), and can be a defensible single-judge fallback for evaluations without access to an ensemble.

\subsection{Epistemic Fragility Beyond the Single-Shot Verdict}
\label{sec:discussion-pressure}

Even though LLM judges are becoming widespread, including being increasingly used as training signals, they are typically deployed as a traditional classifier --- emitting a verdict absent epistemic tests. Finding 2 (FO2) shows that introducing epistemic challenges that actually succeed in changing the judge's mind are almost always likely to be more corruptive than corrective. If L4--L6 is any approximation of multi-agent judging designs like debate-based scalable oversight or self-critique reward loops, our work suggests that there are no LLMs that are epistemically robust or corrective, despite testing on canonical safety judging tasks.

